\section{Introduction}

Direct Preference Optimization (DPO) has rapidly become the dominant paradigm for aligning large language models with human preferences \cite{rafailov2023dpo}. By eliminating the need for a separate reward model and online sampling, DPO reduces alignment to a simple classification-style loss on preference pairs, enabling practitioners to fine-tune models with minimal infrastructure overhead. This simplicity has driven widespread adoption: as of early 2025, the majority of open-weight instruction-tuned models on public repositories employ DPO or its variants as a post-training step \cite{pal2024smaug}. Yet this ubiquity may conceal a latent vulnerability. DPO's theoretical formulation implicitly encourages the policy to concentrate probability mass on a narrow set of preferred completions, a phenomenon known as mode collapse \cite{razin2024vanishing}. While prior work has characterized this collapse as a performance limitation---producing repetitive or degenerate outputs---we argue it constitutes something far more consequential: an exploitable security vulnerability.

Mode collapse under DPO is well-documented. As training progresses, the model's output distribution sharpens around a shrinking set of modes, reducing the effective entropy of its generation process \cite{razin2024vanishing, pal2024smaug}. Standard practice mitigates this through early stopping calibrated to downstream benchmark performance, halting training when helpfulness metrics peak. However, this stopping criterion is agnostic to the distributional properties that govern adversarial robustness. We observe that as output entropy decreases, model behavior becomes increasingly deterministic and therefore increasingly predictable to an adversary. A model with low output entropy conditioned on a given prompt prefix admits a small, enumerable set of likely continuations---precisely the condition that enables targeted exploitation. This connection between alignment-induced collapse and adversarial predictability has not been examined in the security literature.

Concurrent work on LLM poisoning has demonstrated that preference data is a potent attack surface. Pathmanathan et al.\ \cite{pathmanathan2025dpo_poison} showed that contaminating as little as 0.5\% of DPO training pairs suffices to inject targeted misbehaviors, while Wang et al.\ \cite{wang2024rlhfpoison} demonstrated reward model manipulation through preference poisoning. Shadow alignment attacks achieve safety degradation with as few as 100 adversarial samples \cite{yang2023shadow}. Separately, jailbreak methods such as GCG \cite{zou2023gcg} exploit gradient-accessible models to find adversarial suffixes that elicit harmful outputs. What remains unexplored is the intersection of these two research threads: how DPO's intrinsic mode collapse amplifies poisoning efficacy and how the resulting determinism enables a new class of attacks that neither the poisoning nor the jailbreak literature addresses in isolation.

In this paper, we formalize this intersection through a new metric, \textit{Preference Collapse Exploitability} (PCE), which quantifies the security risk arising from entropy reduction during DPO training. Our contributions are fivefold. First, we define PCE and establish its theoretical relationship to adversarial success probability under deterministic output regimes. Second, we empirically demonstrate that PCE increases monotonically with DPO training steps and identify a critical transition point that precedes standard performance peaks---meaning models are routinely trained past the safety threshold. Third, we design \textit{collapse accelerator} attacks: a data poisoning strategy requiring only 1--2\% of training data that actively steers mode collapse toward attacker-specified output modes, achieving targeted harmful generation rates exceeding 85\%. Fourth, we audit six widely-deployed open-source DPO-trained models and find that all exhibit PCE scores above 0.5, placing them in the high-risk regime. Fifth, we propose an entropy-regularized DPO objective that reduces PCE by 59\% while incurring less than 2\% degradation on standard alignment benchmarks, offering a practical defense compatible with existing training pipelines.

The remainder of this paper is organized as follows. Section~\ref{sec:background} reviews DPO training dynamics and threat models for preference learning. Section~\ref{sec:pce} formalizes the PCE metric and its theoretical properties. Section~\ref{sec:collapse_attack} presents the collapse accelerator attack methodology. Section~\ref{sec:experiments} reports empirical results including the open-model audit. Section~\ref{sec:defense} introduces and evaluates our entropy-regularized defense. Section~\ref{sec:discussion} discusses limitations and broader implications, and Section~\ref{sec:conclusion} concludes.
